\chapter{Hoàn thiện yêu cầu mentor và các tính năng tàu}
\label{chap:mentor-final}
\section{Phạm vi và quyền sở hữu phần việc}
Phần tàu được mô tả trong báo cáo gồm điều khiển người chơi, hai mô hình tàu trực quan, ba cấu hình thân tàu, lực đẩy và bánh lái vật lý, nổi trên sóng, tuần tra A/B, A* địa lý, chiến đấu AI, đạn, sát thương, rút lui, cháy và chìm. Các chương trước trình bày cách sử dụng từng hệ thống; chương này bổ sung kết quả tích hợp và yêu cầu mentor. Danh sách là các tính năng hiện có của sản phẩm, không phải tuyên bố một người viết toàn bộ repository. Giao diện Match Settings, kill feed và màn hình tổng kết kế thừa phần việc UI của đồng đội đã vào main; phần tích hợp giữ thiết kế đó và sửa các điểm nối với gameplay.

\section{Đối chiếu yêu cầu mentor}
\begin{longtable}{p{.27\linewidth}p{.65\linewidth}}
\toprule Yêu cầu & Cách xử lý và cách kiểm tra \\ \midrule
Tàu có độ đầm nhưng vẫn dập dềnh & Giảm cản góc thụ động theo từng profile; giữ lực nổi, heave, quán tính, lực đẩy, bánh lái và va chạm. Đo idle, cruise, turn, reverse và reverse-fire trên mặt nước thật. \\
Bỏ cầu debug đỏ khi trúng đạn & Ngắt riêng luồng thực thi DrawDebugSphere trong hai Blueprint đạn có sẵn; nối lại luồng thực thi tiếp theo để giữ sát thương và hiệu ứng. Biên dịch Blueprint và kiểm tra phát bắn thật. \\
Thông báo tàu bị hủy & Dùng kill feed của đồng đội, hiển thị tên bên tiêu diệt và tàu bị phá hủy; trạng thái sống/chết và bộ đếm vẫn do battle subsystem quyết định. \\
Thắng/thua và chơi lại/thoát & Dùng Match Summary, Return to Map Setting, Restart và Quit. Sửa tranh chấp input giữa loading và result; bỏ loading watcher trước khi mở summary. \\
Khôi phục UI & Source hiện có tiêu đề MATCH SETTINGS. Ảnh MATCH SETUP cũ xuất hiện khi mở worktree gốc với DLL cũ. Cần build đúng project, mở đúng worktree và xác minh UI thật trong PIE; không tạo menu thứ hai. \\
\bottomrule
\end{longtable}

\begin{figure}[p]
\centering
\includegraphics[width=.96\linewidth]{../Docs/Integration/images/match-settings.png}
\caption{Match Settings của đồng đội trong PIE ở giai đoạn tích hợp. Ảnh minh họa UI; không thay thế bằng chứng trận đánh tự nhiên.}
\end{figure}
\clearpage
\section{Danh mục tính năng tàu và hướng dẫn demo}
\begin{longtable}{p{.24\linewidth}p{.68\linewidth}}
\toprule Tính năng & Sử dụng và hành vi mong đợi \\ \midrule
Tàu người chơi & Mặc định Medium. Chọn Small ở ô Ship 1 của đội người chơi trước Start New Match để đổi mô hình. Input thực sử dụng Enhanced Input; xem bảng phím ở chương điều khiển. \\
Cấu hình Small/Medium/Large & Chọn từng tàu trong Match Settings. Profile riêng chứa mass, inertia, damping, propulsion, rudder, tốc độ, pontoon, HP và damage. Large/Medium/Small là ba cấu hình gameplay; không đồng nghĩa có ba mô hình nghệ thuật khác nhau. \\
Lực đẩy và bánh lái & Tàu di chuyển bằng lực vật lý hiện hữu. Chỉnh Cruise Speed để đổi tốc độ AI mong muốn; Forward Max Speed là giới hạn, không tự tạo thêm lực. Propulsion Scale và Rudder Scale thay đổi đáp ứng; kiểm tra lại phanh và quay khi chỉnh. \\
Nổi trên sóng & Dùng WaterBodyOcean và các pontoon. Các profile ở \path{/Game/NavalBattle/Profiles}. Cản roll/pitch là tốc độ suy giảm góc, không phải biên độ góc đặt trước. \\
Mô hình Small & Dùng Blueprint người chơi/AI riêng và cấu hình thiết bị phù hợp mô hình Small; giữ thân va chạm, hướng đi và vị trí nòng súng đã tích hợp. Không giả định mô hình static mesh có skeleton. \\
Tuần tra do tác giả đặt & Dừng PIE, đặt tàu và hai TargetPoint khác nhau, gán Patrol Actor A/B trong Navigation. Chọn GENERATE MANUALLY (A-B). Tàu đi từ vị trí đặt đến A, rồi B, rồi A lặp lại. Không đặt hai vùng arrival chồng nhau hoặc ép tàu quay trong khe quá hẹp. \\
Full AI trên map tự đặt & Cần thêm đúng một marker có tag Naval.PlayerPatrol.A và một marker khác có tag Naval.PlayerPatrol.B cho tàu thứ sáu. PlayerStart xác định vị trí xuất phát của tàu này. \\
A* địa lý & Route dựa trên chart biển, hệ ECEF/ENU và kiểm tra nước/va chạm. Di chuyển marker trong PIE phải kích hoạt replanning; di chuyển tàu tới nước hợp lệ phải tái khởi tạo theo vị trí mới. AI phát lệnh qua một pipeline throttle/rudder. \\
Sinh theo seed & Chọn GENERATE FROM SEED, nhập số không âm. Cùng seed và cấu hình đội cho cùng layout; 0 lấy seed mới. Điều đó không bảo đảm mọi va chạm hay kết quả trận đánh giống từng bit. \\
Đội và quy mô & Dùng + và X để đặt 1--3 tàu mỗi đội; chọn đội người chơi hoặc Full AI/Spectator. Không để camera spectator chiếm quyền điều khiển của AI. \\
Chiến đấu AI & Phát hiện, kiểm tra đội/LOS, tính đón đầu đạn đạo, quay súng trong giới hạn cơ khí và bắn bằng hệ đạn chung. Mục tiêu sau đảo không được bắn xuyên địa hình. \\
HP và damage & Standard demo dùng 100 HP / 10 damage; chế độ class dùng thông số profile. Sát thương theo vùng, lọc tự bắn/đồng đội và chống xử lý trùng được kiểm tra riêng. \\
Băng đạn và nạp lại & Theo dõi ammo/reload trên HUD. Standard demo có băng 10 viên và 30 giây nạp lại. Chết phải hủy nạp và chỉ cập nhật alive counter một lần. \\
Rút lui và lùi có bắn & AI chọn ý định rút lui rồi dùng navigation và lực vật lý hiện hữu. Khi an toàn, tiếp tục tuần tra. Bài thử riêng xác nhận phát bắn trong lúc vận tốc tiến âm. \\
Chết, cháy và chìm & Hết HP hoặc capsize hợp lệ kích hoạt vòng đời chết. Có hiệu ứng, chìm và dọn actor; không dùng một khung hình biến mất làm bằng chứng toàn bộ vòng đời. \\
HUD, minimap, spectator & I mở thông tin tàu, Q đổi góc theo đội và Tab chuyển tàu theo logic camera hiện có. Kiểm tra tên, HP, ammo và tàu đang theo dõi khớp nhau; AIController phải được giữ nguyên. \\
Real-Time/Turn-Based & Chọn mode trước Start. Lượt thử 120 giây của cả Player+AI và Full AI ghi nhận 12 lượt luân phiên; tàu ngoài lượt không phát lực đẩy/bánh lái. \\
Menu và kết quả & P mở settings. Back to Game trở lại trận hiện tại; Start New Match tạo trận mới. Sau kết quả dùng Return to Map Setting, Restart hoặc Quit. Quit trong PIE kết thúc PIE, không đóng toàn bộ Editor. \\
\bottomrule
\end{longtable}

\section{Số đo dập dềnh mới nhất}
Small dùng roll damping 1.5/s, pitch damping 0.8/s; Medium 4/s và 5/s; Large 6/s và 4/s. Heave damping vẫn 20/s và lateral damping 4/s. Không khóa yaw, không gán góc sin, không ép tàu luôn nghiêng đúng 2 hoặc 3 độ.
\begin{center}
\begin{tabular}{lrrrr}
\toprule Small & Roll RMS & Pitch RMS & Roll P95 & Pitch P95 \\
\midrule
Idle & 0.069 & 1.117 & 0.125 & 2.034 \\
Cruise & 0.053 & 0.706 & 0.075 & 0.970 \\
Turn & 0.354 & 0.569 & 0.597 & 0.969 \\
Reverse & 0.635 & 1.205 & 0.737 & 2.507 \\
Reverse-fire & 0.653 & 1.179 & 0.793 & 2.374 \\
\bottomrule
\end{tabular}
\end{center}
Đơn vị góc là độ; đây là thống kê góc tuyệt đối, có cả trim trung bình. Peak pitch idle là 2.18 độ và reverse 2.67 độ. Chuyển động phụ thuộc hướng tàu, pha sóng và tốc độ; không phải mọi trục đều đạt 2--3 độ ở mọi thời điểm. Small đã qua kiểm tra hai chiều lái, ngắm, bắn, camera và lùi gần bờ khoảng 99.6 m. Nguồn số liệu chi tiết: \path{Docs/Testing/SMALL_WAVE_FOLLOWUP.json} và trace CSV đi kèm.

\section{Bằng chứng navigation, sửa deadlock và giới hạn kết luận}
Lượt chạy map đã lưu với sáu Medium đạt ít nhất bốn arrival A/B mỗi tàu trong 1350.93 giây mô phỏng. Tất cả giữ controller, tham chiếu marker và physics; tỷ lệ mẫu có nước thấp nhất 99.974\%. Peak pitch/roll trong toàn bộ thao tác và đi sát nhau là 3.05/4.49 độ, khác với mức điển hình khi idle. Chỉ TP\_Blue02\_B được dời khoảng 645 m để cách marker A đối phương 481 m; vị trí xuất phát của các tàu và các marker khác được giữ.

Hai lỗi độc lập được tách khỏi tuning sóng: tàu cùng hướng vẫn nhường tàu đứng yên dù có traffic priority; và escort dùng hết ba lần thử route rồi ngừng vĩnh viễn. Lỗi đầu được sửa bằng điều kiện lead vessel thực sự đang chạy, vẫn giữ hull sweep và emergency brake. Lỗi sau mở lại cửa sổ thử sau cooldown bằng planner có sẵn, giữ quyền điều khiển của combat/search đang hoạt động. Bài thử hàng đợi vượt 100 m và bài thử retry vượt 50 m đều dùng lực vật lý thật.

Trận 3v3 tự nhiên trên DLL cuối với seed 20261028 đã kết thúc BLUE WINS sau 803.00 giây mô phỏng. Cả sáu tàu bắn; ba tàu Red bị tiêu diệt do hết máu, không do capsize hay sát thương do bài thử chèn vào. Controller và physics được giữ trong suốt thời gian sống. Summary chỉ xuất hiện một lần, bộ đếm Blue 3 / Red 0 chính xác; thao tác chuột Return to Map Setting mở lại đúng Match Settings.\par

\textbf{Phân biệt loại bằng chứng:} build C++, compile Blueprint, automation cô lập, input PIE, fault injection và trận đánh tự nhiên là các mức khác nhau. Bài thử fault chủ động gây lỗi hoặc sát thương không được dùng để tuyên bố trận thắng tự nhiên. Trận tự nhiên bị dừng hoặc bị thay world từ Editor phải ghi interrupted. Nghiệm thu cuối và trạng thái Git được cập nhật trong \path{Docs/Testing/UI_BUOYANCY_REGRESSION.md} và \path{Docs/Integration/GIT_MERGE_REPORT.md}; không coi source cleanup hoàn tất chỉ vì đã dọn Tools/Docs.

\section{Quy trình demo khuyến nghị}
\begin{enumerate}
\item Mở đúng project/worktree có DLL mới, chạy MainMapLevel và kiểm tra tiêu đề MATCH SETTINGS cùng ba nhóm General, Deploy Troops, Team Setup.
\item Chọn Medium mặc định hoặc Small ở slot người chơi; chọn Player+AI để demo lái, ngắm và bắn. Dùng I để hiện thông tin tàu.
\item Chọn Manual để demo tuyến A/B do tác giả đặt; dùng Full AI để quan sát toàn bộ sáu tàu và đổi camera mà không chiếm quyền AI.
\item Chọn seed đã được ghi nhận cho demo giao chiến. Quan sát contact, phát bắn, reload, rút lui, kill feed, cháy/chìm và bộ đếm đội.
\item Sau khi một đội hết tàu, kiểm tra summary và thao tác Return/Restart/Quit. Giữ nguyên lượt kiểm tra tự động đang chạy; đổi hoặc dừng PIE giữa chừng làm bằng chứng đó bị gián đoạn.
\end{enumerate}

\clearpage
\section{Tích hợp và sử dụng source}
UI nằm trong \path{ShipSetupWidget.cpp}, được \path{PlayerController} tạo; lựa chọn đi qua \path{NavalMatchSettings}, game mode, roster và battlefield subsystem. Physics profile thuộc \path{ShipPhysicsProfile}; lực và cản thuộc \path{ShipMovementComponent}. Route/steering thuộc \path{ShipEnemyNavigationComponent} và \path{SeaNavigationMap}; ý định đội thuộc \path{NavalBattlefieldTactics}. Súng và sát thương thuộc \path{ShipCombatComponent}, \path{ShipWeaponComponent} và \path{ShipProjectileDamageComponent}. Tách các trách nhiệm này giúp giữ một đường phát lệnh di chuyển.

Phải commit thay đổi đã kiểm tra, fetch và merge main vào nhánh feature, giải quyết conflict trên feature, rồi build/PIE lại trước PR vào main. Không text-merge uasset/umap, không reset/clean để bỏ map, không force-push. Font local được sao lưu và bảo toàn; DLL cũ không được dùng làm bằng chứng UI source đã bị mất. Báo cáo không chứng nhận cooked build, multiplayer, mọi seed hay mọi cấu hình sóng ngoài các lượt thử đã ghi.
